\begin{tikzpicture}[x={\dimexpr\linewidth/171\relax},y=1mm,
  font=\figurefont\footnotesize,text=NNInk,
  upsample/.style={draw=NNInput,fill=NNInput!18,line width=1.1pt,
    rounded corners=1mm,minimum width=12mm,minimum height=9mm,
    inner sep=1pt,align=center}]
  \cnnArchitecture{54}{7}
  \cnnFeature{in}{2}{17}{6}{28}{1}{NNInput}{输入}{$64\times64$\\$c_{\mathrm{in}}$通道}
  \cnnFeature{e0}{13}{17}{9}{28}{2}{NNHidden}{编码1}{$64\times64$\\$32$通道}
  \cnnFeature{e1}{30}{21}{8}{20}{3}{NNHidden}{编码2}{$32\times32$\\$64$通道}
  \cnnFeature{e2}{46}{24.5}{7}{13}{4}{NNHidden}{瓶颈}{$16\times16$\\$128$通道}
  \cnnFeature{d1}{92}{21}{8}{20}{3}{NNHidden}{解码2}{$32\times32$\\$64$通道}
  \cnnFeature{d0}{139}{17}{9}{28}{2}{NNHidden}{解码1}{$64\times64$\\$32$通道}
  \cnnFeature{out}{156.5}{17}{6}{28}{2}{NNOutput}{输出}{$64\times64$\\$q$或2通道}
  \foreach \a/\b in {in/e0,e0/e1,e1/e2,d0/out}{
    \cnnScaleLinks{\a}{\b}\cnnMap{\a}{\b}
  }
  % 每级解码都显式经过上采样、同尺度拼接和卷积融合。
  \node[upsample] (up1) at (65,31) {上采样};
  \node[upsample] (up0) at (112,31) {上采样};
  \nnConcatOp[6mm]{cat1}{79}{31}
  \nnConcatOp[6mm]{cat0}{126}{31}
  \draw[nn flow] (e2-east)--(up1.west);
  \draw[nn flow] (up1.east)--(cat1.west);
  \draw[nn flow,draw=NNHidden] (cat1.east)--(d1-west)
    node[midway,above=1.2mm,cnn operation] {卷积};
  \draw[nn flow] (d1-east)--(up0.west);
  \draw[nn flow] (up0.east)--(cat0.west);
  \draw[nn flow,draw=NNHidden] (cat0.east)--(d0-west)
    node[midway,above=1.2mm,cnn operation] {卷积};
  % 编码特征从张量上方取出，再沿顶部路由带进入拼接节点。
  \draw[cnn skip] (e1-north)--(35.5,62)--(79,62)--(cat1.north);
  \draw[cnn skip] (e0-north)--(18.5,70)--(126,70)--(cat0.north);
  \node[cnn operation] at (153,14) {$1\times1$};
\end{tikzpicture}
